\documentclass[11pt]{article}

\usepackage{enumitem}
\usepackage{longtable}
\usepackage{booktabs}
\usepackage{hyperref}
\usepackage[margin=1in]{geometry}

\title{Algebraic Geometry II\\
\large Course Design and Topic Breakdown}
\author{}
\date{}

\begin{document}

\maketitle

\section*{Course Overview}

\textbf{Course level:} PhD

\textbf{Format:} Lecture course only; two 75-minute lectures per week.

\textbf{Duration:} 15 weeks, 30 lectures. The syllabus indicates 14--15 weeks; this design uses 15 weeks to provide a complete two-lecture-per-week structure.

\textbf{Primary text:} Robin Hartshorne, \emph{Algebraic Geometry}, primarily Chapters II and III.

\textbf{Assessment:} Weekly problem sets and an oral final examination. No written midterm or separate project is assumed.

\textbf{Prerequisites:} Students are expected to have a working knowledge of commutative algebra, point-set topology, and classical algebraic geometry. Familiarity with the basics of category theory is useful.

\section*{Course Goals}

The course introduces the modern language of schemes and the theory of cohomology in algebraic geometry, following the development associated with Grothendieck. By the end of the course, students should be able to work fluently with schemes, sheaves, and morphisms of schemes; translate geometric questions into algebraic statements using local and global methods; and use sheaf cohomology to formulate and prove fundamental results about algebraic varieties and schemes.

\section*{Assumptions Used in This Design}

\begin{itemize}
    \item The 14--15 week range is resolved as a 15-week semester with 30 lectures.
    \item The course follows Hartshorne Chapters II and III as its primary organizational framework, with selected topics emphasized according to their conceptual importance rather than attempting to cover every result in those chapters at equal depth.
    \item Weekly problem sets are assumed to be assigned throughout the semester and are used primarily for technical mastery and proof practice.
    \item The oral final examination is assumed to test conceptual understanding, the ability to state and explain major definitions and theorems, and the ability to give short proofs or proof sketches.
    \item The course assumes substantial prior commutative algebra, including localization, tensor products, integral extensions, Noetherian rings, modules, and basic dimension theory.
    \item Classical algebraic geometry is assumed at the level of affine/projective varieties, regular functions, morphisms, tangent-space intuition, and basic algebra--geometry correspondences.
\end{itemize}

\section{Module 1: From Classical Algebraic Geometry to Schemes}

\textbf{Weeks:} 1--2

\textbf{Goal:} Transition from varieties to the local and categorical language of schemes by introducing sheaves, locally ringed spaces, and the spectrum of a ring.

\subsection{Topological Spaces and Sheaves}

\textbf{Core concepts}
\begin{itemize}
    \item Topological spaces and local data.
    \item Presheaves and sheaves of rings and modules.
    \item Stalks and germs.
    \item Sheafification.
    \item Sections over open subsets.
    \item Restriction maps and gluing.
\end{itemize}

\textbf{Learning objectives}

Students should be able to:
\begin{itemize}
    \item Define a presheaf and a sheaf and explain the difference between them.
    \item Compute stalks in basic examples.
    \item Construct and use the sheafification of a presheaf.
    \item Verify the sheaf axioms in concrete examples.
    \item Explain why sheaves are appropriate for encoding locally defined algebraic functions.
\end{itemize}

\textbf{Prerequisites:} Point-set topology and basic ring/module theory.

\textbf{Difficulty analysis:} Students often confuse a sheaf with its collection of global sections and underestimate the role of stalks. The abstraction also introduces several layers of structure simultaneously.

\textbf{Mitigation:} Begin with familiar sheaves of continuous functions and regular functions before introducing sheaves of rings abstractly. Repeatedly distinguish global sections, sections over an open set, and stalks.

\textbf{Assessment:} Week 1 problem set. Exercises should include stalk calculations, sheaf verification, and construction of examples.

\textbf{Example questions}
\begin{enumerate}
    \item Explain why the assignment $U \mapsto C(U,\mathbb{R})$ defines a sheaf. Identify its stalk at a point.
    \item Give an example of a presheaf that fails the sheaf axiom and identify precisely where the failure occurs.
\end{enumerate}

\subsection{Affine Schemes}

\textbf{Core concepts}
\begin{itemize}
    \item Prime spectrum $\operatorname{Spec}(A)$.
    \item Zariski topology.
    \item Distinguished open subsets $D(f)$.
    \item Structure sheaf $\mathcal{O}_{\operatorname{Spec}(A)}$.
    \item Locally ringed spaces.
    \item The fundamental correspondence between rings and affine schemes.
\end{itemize}

\textbf{Learning objectives}

Students should be able to:
\begin{itemize}
    \item Define $\operatorname{Spec}(A)$ and describe its Zariski topology.
    \item Compute distinguished open subsets and their basic properties.
    \item Describe the structure sheaf on affine schemes.
    \item Prove or explain why the stalk at $\mathfrak p$ is $A_{\mathfrak p}$.
    \item Use the equivalence between commutative rings and affine schemes in concrete arguments.
\end{itemize}

\textbf{Prerequisites:} Localization, prime ideals, maximal ideals, Zariski topology, and sheaves.

\textbf{Difficulty analysis:} The topology of $\operatorname{Spec}(A)$ behaves very differently from familiar Hausdorff spaces, and students may initially regard prime ideals merely as points rather than geometric objects carrying local algebra.

\textbf{Mitigation:} Work through $\operatorname{Spec}(\mathbb Z)$, polynomial rings, local rings, and finite rings before treating general statements.

\textbf{Assessment:} Week 2 problem set. Focus on explicit spectra, distinguished opens, and stalk computations.

\textbf{Example questions}
\begin{enumerate}
    \item Describe the points and topology of $\operatorname{Spec}(\mathbb Z)$.
    \item Show that $D(f)$ is naturally identified with $\operatorname{Spec}(A_f)$.
    \item Explain the geometric meaning of localization at a prime ideal.
\end{enumerate}

\section{Module 2: Schemes and Their Morphisms}

\textbf{Weeks:} 3--4

\textbf{Goal:} Develop schemes as locally ringed spaces and understand morphisms, closed subschemes, and the relationship between algebraic and geometric constructions.

\subsection{Schemes as Locally Ringed Spaces}

\textbf{Core concepts}
\begin{itemize}
    \item Locally ringed spaces.
    \item Schemes.
    \item Open and closed subschemes.
    \item Reduced and irreducible schemes.
    \item Generic points.
    \item Nilpotents and nonreduced schemes.
\end{itemize}

\textbf{Learning objectives}

Students should be able to:
\begin{itemize}
    \item Define a scheme as a locally ringed space locally isomorphic to an affine scheme.
    \item Construct basic examples of schemes.
    \item Distinguish the underlying topological space from its structure sheaf.
    \item Explain the geometric significance of nilpotents.
    \item Identify reduced, irreducible, and integral schemes.
\end{itemize}

\textbf{Prerequisites:} Affine schemes and commutative algebra.

\textbf{Difficulty analysis:} The same topological space can support different scheme structures, which is conceptually difficult for students coming from classical varieties.

\textbf{Mitigation:} Compare $\operatorname{Spec}(k[\epsilon]/(\epsilon^2))$ with $\operatorname{Spec}(k)$ and explicitly separate topological and sheaf-theoretic information.

\textbf{Assessment:} Week 3 problem set.

\textbf{Example questions}
\begin{enumerate}
    \item Compare $\operatorname{Spec}(k)$ and $\operatorname{Spec}(k[\epsilon]/(\epsilon^2))$ as schemes.
    \item Explain why a scheme may contain geometric information invisible in its underlying topological space.
\end{enumerate}

\subsection{Morphisms of Schemes}

\textbf{Core concepts}
\begin{itemize}
    \item Morphisms of locally ringed spaces.
    \item Morphisms of schemes.
    \item Morphisms of affine schemes and ring homomorphisms.
    \item Global sections.
    \item Functoriality of $\operatorname{Spec}$.
    \item Closed immersions and affine morphisms.
\end{itemize}

\textbf{Learning objectives}

Students should be able to:
\begin{itemize}
    \item Define a morphism of schemes.
    \item Construct the morphism induced by a ring homomorphism.
    \item Translate affine geometric statements into algebraic statements.
    \item Determine whether a given map of affine schemes is a closed immersion.
    \item Explain the contravariant relationship between rings and affine schemes.
\end{itemize}

\textbf{Prerequisites:} Locally ringed spaces and affine schemes.

\textbf{Difficulty analysis:} Morphisms contain both a continuous map and a morphism of sheaves, and students often focus only on the underlying point map.

\textbf{Mitigation:} Use affine examples first and explicitly track both components of a morphism.

\textbf{Assessment:} Week 4 problem set.

\textbf{Example questions}
\begin{enumerate}
    \item Given $\varphi:A\to B$, construct the associated morphism $\operatorname{Spec}(B)\to\operatorname{Spec}(A)$.
    \item Determine which ring maps correspond to closed immersions of affine schemes.
\end{enumerate}

\section{Module 3: Geometric Constructions with Schemes}

\textbf{Weeks: 5--6}

\textbf{Goal:} Develop fiber products, base change, and geometric properties of morphisms, establishing the categorical flexibility of schemes.

\subsection{Fiber Products}

\textbf{Core concepts}
\begin{itemize}
    \item Products of schemes.
    \item Fiber products.
    \item Tensor products and affine fiber products.
    \item Base change.
    \item Fibers of morphisms.
    \item Universal properties.
\end{itemize}

\textbf{Learning objectives}

Students should be able to:
\begin{itemize}
    \item Define a fiber product using its universal property.
    \item Compute fiber products of affine schemes.
    \item Relate tensor products of rings to geometric base change.
    \item Compute fibers of morphisms.
    \item Use universal properties to establish uniqueness and functoriality.
\end{itemize}

\textbf{Prerequisites:} Tensor products, affine schemes, and basic category theory.

\textbf{Difficulty analysis:} Universal properties and tensor products require students to simultaneously reason categorically and algebraically.

\textbf{Mitigation:} Alternate abstract universal-property arguments with explicit affine calculations.

\textbf{Assessment:} Week 5 problem set.

\textbf{Example questions}
\begin{enumerate}
    \item Compute $\operatorname{Spec}(B)\times_{\operatorname{Spec}(A)}\operatorname{Spec}(C)$.
    \item Explain how the fiber of a morphism over a point is obtained by base change.
    \item Construct a fiber product in a simple example and verify its universal property.
\end{enumerate}

\subsection{Properties of Morphisms}

\textbf{Core concepts}
\begin{itemize}
    \item Finite type and finite morphisms.
    \item Separated morphisms.
    \item Proper morphisms.
    \item Affine morphisms.
    \item Closedness and valuative intuition.
    \item Stability under composition and base change.
\end{itemize}

\textbf{Learning objectives}

Students should be able to:
\begin{itemize}
    \item Define finite type, finite, affine, separated, and proper morphisms.
    \item Compare these properties and identify implications between them.
    \item Determine these properties in standard affine examples.
    \item Explain why properness is a scheme-theoretic analogue of compactness.
    \item Use stability properties under composition and base change.
\end{itemize}

\textbf{Prerequisites:} Fiber products, Noetherian rings, integral extensions, and classical projective geometry.

\textbf{Difficulty analysis:} The definitions are numerous and superficially similar, while their geometric interpretations are subtle.

\textbf{Mitigation:} Maintain a running implication diagram and pair every formal definition with an explicit geometric example and counterexample.

\textbf{Assessment:} Week 6 problem set.

\textbf{Example questions}
\begin{enumerate}
    \item Explain the relationship between finite and finite-type morphisms.
    \item Give a geometric interpretation of separatedness.
    \item Why is properness naturally compared with compactness?
\end{enumerate}

\section{Module 4: Projective and Quasi-Projective Schemes}

\textbf{Weeks: 7--8}

\textbf{Goal:} Connect the abstract scheme-theoretic language to projective geometry and develop the geometric properties needed for cohomological applications.

\subsection{Projective Schemes and Proj}

\textbf{Core concepts}
\begin{itemize}
    \item Graded rings.
    \item Homogeneous prime ideals.
    \item $\operatorname{Proj}(S)$.
    \item Structure sheaf on $\operatorname{Proj}(S)$.
    \item Projective space.
    \item Twisting sheaves $\mathcal{O}(n)$.
\end{itemize}

\textbf{Learning objectives}

Students should be able to:
\begin{itemize}
    \item Define $\operatorname{Proj}(S)$.
    \item Describe its standard affine open subsets.
    \item Construct projective space scheme-theoretically.
    \item Explain the role of grading in projective geometry.
    \item Work with the twisting sheaves $\mathcal{O}(n)$.
\end{itemize}

\textbf{Prerequisites:} Graded rings, localization, affine schemes, and classical projective geometry.

\textbf{Difficulty analysis:} $\operatorname{Proj}$ introduces homogeneous localization and irrelevant ideals while requiring students to translate between graded algebra and geometry.

\textbf{Mitigation:} Construct $\mathbb{P}^n$ explicitly and repeatedly compare $\operatorname{Spec}$ and $\operatorname{Proj}$.

\textbf{Assessment:} Week 7 problem set.

\textbf{Example questions}
\begin{enumerate}
    \item Construct $\mathbb{P}^n_k$ as $\operatorname{Proj}(k[x_0,\ldots,x_n])$.
    \item Identify the standard affine open subsets of $\operatorname{Proj}(S)$.
    \item Explain why ordinary localization is insufficient to describe projective geometry directly.
\end{enumerate}

\subsection{Projective Morphisms and Classical Correspondence}

\textbf{Core concepts}
\begin{itemize}
    \item Projective morphisms.
    \item Closed subschemes of projective space.
    \item Quasi-projective schemes.
    \item Very ample line bundles.
    \item Relationship between projective varieties and projective schemes.
\end{itemize}

\textbf{Learning objectives}

Students should be able to:
\begin{itemize}
    \item Define projective and quasi-projective schemes.
    \item Recognize standard projective constructions.
    \item Explain how homogeneous ideals determine closed subschemes of projective space.
    \item Relate projective embeddings to ample and very ample line bundles.
\end{itemize}

\textbf{Prerequisites:} $\operatorname{Proj}$ and classical projective algebraic geometry.

\textbf{Difficulty analysis:} Students must coordinate several equivalent viewpoints: embeddings, graded rings, homogeneous ideals, and invertible sheaves.

\textbf{Mitigation:} Use one running example, such as a projective curve, through all four viewpoints.

\textbf{Assessment:} Week 8 problem set.

\textbf{Example questions}
\begin{enumerate}
    \item Describe the closed subscheme of $\mathbb{P}^n$ associated with a homogeneous ideal.
    \item Explain how an embedding into projective space produces a very ample line bundle.
\end{enumerate}

\section{Module 5: Sheaves of Modules and Quasi-Coherent Sheaves}

\textbf{Weeks: 9--10}

\textbf{Goal:} Develop quasi-coherent and coherent sheaves as the natural algebraic objects on schemes and establish their affine correspondence with modules.

\subsection{Sheaves of Modules and Quasi-Coherent Sheaves}

\textbf{Core concepts}
\begin{itemize}
    \item $\mathcal{O}_X$-modules.
    \item Quasi-coherent sheaves.
    \item Coherent sheaves.
    \item The sheaf associated to a module.
    \item Localization of modules.
    \item Exactness and local properties.
\end{itemize}

\textbf{Learning objectives}

Students should be able to:
\begin{itemize}
    \item Define $\mathcal{O}_X$-modules, quasi-coherent sheaves, and coherent sheaves.
    \item Construct the sheaf $\widetilde{M}$ associated to an $A$-module $M$.
    \item Prove or use the correspondence between modules and quasi-coherent sheaves on affine schemes.
    \item Determine whether common sheaves are quasi-coherent or coherent.
    \item Translate module-theoretic constructions into sheaf-theoretic ones.
\end{itemize}

\textbf{Prerequisites:} Modules, localization, sheaves, and affine schemes.

\textbf{Difficulty analysis:} The distinction between arbitrary sheaves of modules and quasi-coherent sheaves can seem artificial until students see its compatibility with localization.

\textbf{Mitigation:} Establish the affine correspondence concretely and use it as the central organizing principle.

\textbf{Assessment:} Week 9 problem set.

\textbf{Example questions}
\begin{enumerate}
    \item Given an $A$-module $M$, describe $\widetilde{M}$ on $D(f)$.
    \item Show that quasi-coherence is a local condition.
    \item Explain why quasi-coherent sheaves are the correct scheme-theoretic analogue of algebraic vector-valued functions.
\end{enumerate}

\subsection{Exact Sequences and Coherent Sheaves}

\textbf{Core concepts}
\begin{itemize}
    \item Exact sequences of sheaves.
    \item Kernels, cokernels, and images.
    \item Coherent sheaves on Noetherian schemes.
    \item Locally free sheaves.
    \item Vector bundles.
    \item Locally free resolutions.
\end{itemize}

\textbf{Learning objectives}

Students should be able to:
\begin{itemize}
    \item Determine exactness of a sequence of sheaves locally.
    \item Relate coherent sheaves to finite modules on affine schemes.
    \item Distinguish coherent, locally free, and quasi-coherent sheaves.
    \item Interpret locally free sheaves geometrically as vector bundles.
    \item Construct and use short exact sequences of coherent sheaves.
\end{itemize}

\textbf{Prerequisites:} Quasi-coherent sheaves and module theory.

\textbf{Difficulty analysis:} Students may incorrectly assume that exactness of sheaves can always be checked using global sections.

\textbf{Mitigation:} Emphasize stalkwise exactness and give examples where taking global sections destroys exactness.

\textbf{Assessment:} Week 10 problem set.

\textbf{Example questions}
\begin{enumerate}
    \item Explain why a sequence of sheaves is exact if and only if it is exact on every stalk.
    \item Give an example showing that global sections need not preserve exact sequences.
\end{enumerate}

\section{Module 6: Cohomology of Sheaves}

\textbf{Weeks: 11--12}

\textbf{Goal:} Introduce sheaf cohomology as a systematic measure of the failure of global sections to be exact and develop its basic computational machinery.

\subsection{Derived Functors and Sheaf Cohomology}

\textbf{Core concepts}
\begin{itemize}
    \item Left and right exact functors.
    \item Injective resolutions.
    \item Derived functors.
    \item Cohomology groups $H^i(X,\mathcal{F})$.
    \item Long exact sequences in cohomology.
    \item Functoriality.
\end{itemize}

\textbf{Learning objectives}

Students should be able to:
\begin{itemize}
    \item Explain why sheaf cohomology arises from the global sections functor.
    \item Define $H^i(X,\mathcal{F})$ using an injective resolution.
    \item Derive the long exact sequence associated with a short exact sequence of sheaves.
    \item Explain the functoriality of cohomology.
    \item Interpret $H^1$ and higher groups as obstructions to global constructions.
\end{itemize}

\textbf{Prerequisites:} Basic homological algebra, exact sequences, modules, and quasi-coherent sheaves.

\textbf{Difficulty analysis:} Derived functors introduce a significant abstraction jump and students may focus on the formal resolution rather than understanding what cohomology measures.

\textbf{Mitigation:} Motivate cohomology through the failure of global sections to preserve exactness before introducing resolutions.

\textbf{Assessment:} Week 11 problem set.

\textbf{Example questions}
\begin{enumerate}
    \item Starting from a short exact sequence
    \[
    0\to\mathcal{F}'\to\mathcal{F}\to\mathcal{F}''\to0,
    \]
    explain the construction of the associated long exact sequence.
    \item Why is $H^0(X,\mathcal{F})$ identified with $\Gamma(X,\mathcal{F})$?
    \item Explain what information $H^1(X,\mathcal{F})$ can contain that is invisible to global sections alone.
\end{enumerate}

\subsection{Cech Cohomology and Affine Vanishing}

\textbf{Core concepts}
\begin{itemize}
    \item Cech complexes.
    \item Cohomology associated to affine covers.
    \item Cohomology of quasi-coherent sheaves on affine schemes.
    \item Vanishing of higher cohomology on affine schemes.
    \item Comparison of Cech and sheaf cohomology.
\end{itemize}

\textbf{Learning objectives}

Students should be able to:
\begin{itemize}
    \item Construct the Cech complex associated to an open cover.
    \item Compute low-degree cohomology using Cech methods.
    \item State and apply the affine vanishing theorem.
    \item Explain why affine schemes are cohomologically simple for quasi-coherent sheaves.
    \item Use cohomological vanishing to convert geometric questions into algebraic ones.
\end{itemize}

\textbf{Prerequisites:} Sheaf cohomology, complexes, localization, and quasi-coherent sheaves.

\textbf{Difficulty analysis:} Students must keep track of alternating restriction maps and intersections of open sets while simultaneously interpreting the resulting algebra.

\textbf{Mitigation:} Work through two-open-set covers explicitly before introducing general Cech complexes.

\textbf{Assessment:} Week 12 problem set.

\textbf{Example questions}
\begin{enumerate}
    \item Construct the Cech complex for a cover $X=U\cup V$.
    \item Explain why $H^i(X,\widetilde{M})=0$ for $i>0$ when $X$ is affine.
    \item Compute a low-degree Cech cohomology group for a simple cover.
\end{enumerate}

\section{Module 7: Cohomology of Projective Schemes}

\textbf{Weeks: 13--14}

\textbf{Goal:} Apply sheaf cohomology to projective geometry, compute cohomology of twisting sheaves, and develop the fundamental finiteness and vanishing results of projective geometry.

\subsection{Cohomology of Projective Space}

\textbf{Core concepts}
\begin{itemize}
    \item Cohomology of $\mathcal{O}_{\mathbb{P}^n}(d)$.
    \item Serre's computation of cohomology on projective space.
    \item Hyperplane sections.
    \item Twisting sheaves.
    \item Euler characteristic.
    \item Serre duality in projective space.
\end{itemize}

\textbf{Learning objectives}

Students should be able to:
\begin{itemize}
    \item State the cohomology groups of $\mathcal{O}_{\mathbb{P}^n}(d)$ in the relevant ranges.
    \item Compute spaces of global sections of twisting sheaves.
    \item Use short exact sequences involving hyperplanes to compute cohomology.
    \item Explain the role of twisting in controlling positivity and vanishing.
    \item Apply duality to relate different cohomology groups.
\end{itemize}

\textbf{Prerequisites:} Projective schemes, coherent sheaves, sheaf cohomology, and basic homological algebra.

\textbf{Difficulty analysis:} Cohomology calculations combine combinatorial dimension formulas, exact sequences, and geometric intuition. The pattern of vanishing and nonvanishing can initially appear memorization-based.

\textbf{Mitigation:} Derive the formulas repeatedly from hyperplane exact sequences rather than presenting them only as a table to memorize.

\textbf{Assessment:} Week 13 problem set.

\textbf{Example questions}
\begin{enumerate}
    \item Compute $H^0(\mathbb{P}^n,\mathcal{O}(d))$ for $d\geq 0$ and determine its dimension.
    \item Use a hyperplane exact sequence to relate the cohomology of $\mathcal{O}_{\mathbb{P}^n}(d)$ to that of $\mathcal{O}_{\mathbb{P}^{n-1}}(d)$.
    \item Explain the pattern of vanishing for intermediate cohomology of $\mathcal{O}_{\mathbb{P}^n}(d)$.
\end{enumerate}

\subsection{Serre Vanishing and Finiteness}

\textbf{Core concepts}
\begin{itemize}
    \item Serre vanishing.
    \item Global generation.
    \item Twisting coherent sheaves.
    \item Finite-dimensionality of cohomology.
    \item Hilbert polynomials.
    \item Relation between cohomology and graded modules.
\end{itemize}

\textbf{Learning objectives}

Students should be able to:
\begin{itemize}
    \item State Serre's vanishing theorem.
    \item Explain why twisting eventually eliminates higher cohomology.
    \item Apply Serre vanishing to coherent sheaves on projective schemes.
    \item Relate Hilbert functions and Hilbert polynomials to cohomology.
    \item Explain the geometric significance of finite-dimensional cohomology.
\end{itemize}

\textbf{Prerequisites:} Coherent sheaves, projective schemes, and basic Hilbert polynomial theory.

\textbf{Difficulty analysis:} Students may view Serre vanishing as an isolated technical theorem rather than part of a larger relationship between positivity, global generation, and cohomological simplicity.

\textbf{Mitigation:} Organize the discussion around the question of what sufficiently positive twisting does to a coherent sheaf.

\textbf{Assessment:} Week 14 problem set.

\textbf{Example questions}
\begin{enumerate}
    \item State Serre vanishing precisely and identify all hypotheses.
    \item Explain how sufficiently positive twists affect the global sections and higher cohomology of a coherent sheaf.
    \item Relate the Hilbert polynomial of a coherent sheaf to its Euler characteristic.
\end{enumerate}

\section{Module 8: Synthesis and Advanced Applications}

\textbf{Week:} 15

\textbf{Goal:} Consolidate the scheme-theoretic and cohomological framework and prepare students for the oral final by connecting the major constructions and theorems.

\subsection{Schemes, Cohomology, and the Geometry--Algebra Dictionary}

\textbf{Core concepts}
\begin{itemize}
    \item Affine and projective schemes.
    \item Morphisms and fiber products.
    \item Quasi-coherent and coherent sheaves.
    \item Global sections and cohomology.
    \item Vanishing theorems.
    \item Algebra--geometry correspondences.
    \item Local-to-global reasoning.
\end{itemize}

\textbf{Learning objectives}

Students should be able to:
\begin{itemize}
    \item Synthesize the definitions introduced throughout the course into a coherent scheme-theoretic framework.
    \item Explain how local algebra, global geometry, and cohomology interact.
    \item Select appropriate scheme-theoretic tools for a given geometric problem.
    \item Reconstruct major proofs or proof strategies from the course.
    \item Explain the conceptual relationships among affine schemes, projective schemes, sheaves, morphisms, and cohomology.
\end{itemize}

\textbf{Prerequisites:} All preceding modules.

\textbf{Difficulty analysis:} The principal difficulty is integration rather than a new technical concept. Students may know individual definitions and theorems without seeing how they form a unified language.

\textbf{Mitigation:} Use synthesis problems that require several constructions in sequence and conduct oral-exam-style review sessions.

\textbf{Assessment:} Final weekly problem set plus preparation for the oral final examination.

\textbf{Example questions}
\begin{enumerate}
    \item Explain the chain of ideas that leads from a commutative ring $A$ to the scheme $\operatorname{Spec}(A)$, its structure sheaf, and its quasi-coherent sheaves.
    \item Explain why sheaf cohomology is naturally viewed as a local-to-global theory.
    \item Given a geometric problem on a projective scheme, describe a plausible sequence of scheme-theoretic and cohomological tools for attacking it.
\end{enumerate}

\section{Prerequisite Dependency Map}

\begin{longtable}{p{0.27\textwidth}p{0.38\textwidth}p{0.25\textwidth}}
\toprule
\textbf{Topic} & \textbf{Prerequisite(s)} & \textbf{Where Taught} \\
\midrule
Sheaves & Point-set topology; basic algebra & Assumed background \\
Stalks and germs & Topology; presheaves & Module 1 \\
$\operatorname{Spec}(A)$ & Commutative algebra; prime ideals; localization & Assumed background + Module 1 \\
Structure sheaf & Sheaves; localization; $\operatorname{Spec}(A)$ & Module 1 \\
Locally ringed spaces & Sheaves of rings; local rings & Module 1 \\
Schemes & Locally ringed spaces; affine schemes & Module 2 \\
Morphisms of schemes & Locally ringed spaces; affine schemes & Module 2 \\
Closed immersions & Affine morphisms; quotient rings & Module 2 \\
Fiber products & Tensor products; categorical universal properties & Module 3 \\
Base change & Fiber products; tensor products & Module 3 \\
Finite morphisms & Finite algebras; affine morphisms & Assumed background + Module 3 \\
Separated morphisms & Diagonal morphism; fiber products & Module 3 \\
Proper morphisms & Separatedness; finite type; projective geometry & Module 3 \\
$\operatorname{Proj}$ & Graded rings; localization; projective geometry & Module 4 \\
Projective schemes & $\operatorname{Proj}$; homogeneous ideals & Module 4 \\
Quasi-coherent sheaves & Sheaves of modules; localization & Module 5 \\
Coherent sheaves & Noetherian algebra; quasi-coherent sheaves & Module 5 \\
Locally free sheaves & Modules; localizations; coherent sheaves & Module 5 \\
Sheaf cohomology & Homological algebra; exact sequences & Module 6 \\
Derived functors & Homological algebra & Assumed background + Module 6 \\
Cech cohomology & Sheaves; complexes; cohomology & Module 6 \\
Affine vanishing & Quasi-coherent sheaves; Cech cohomology & Module 6 \\
Cohomology of projective space & $\operatorname{Proj}$; coherent sheaves; cohomology & Module 7 \\
Serre vanishing & Projective schemes; coherent sheaves & Module 7 \\
Hilbert polynomial applications & Graded algebra; projective geometry; cohomology & Modules 4, 5, and 7 \\
Course synthesis & All preceding topics & Module 8 \\
\bottomrule
\end{longtable}

\section{Assessment Plan}

The course uses weekly problem sets and an oral final examination. Because the course is proof-oriented and at the PhD level, the problem sets are intended to develop technical fluency, while the oral examination emphasizes conceptual organization, theorem statements, proof strategies, and the ability to reason aloud.

\begin{longtable}{p{0.15\textwidth}p{0.28\textwidth}p{0.27\textwidth}p{0.20\textwidth}}
\toprule
\textbf{Week} & \textbf{Primary Content} & \textbf{Assessment} & \textbf{Purpose} \\
\midrule
1 & Sheaves and stalks & Problem Set 1 & Establish facility with local-to-global language. \\
2 & Affine schemes & Problem Set 2 & Test explicit spectrum and structure-sheaf calculations. \\
3 & Schemes & Problem Set 3 & Test understanding of scheme structures and local properties. \\
4 & Morphisms & Problem Set 4 & Test translation between ring maps and geometric maps. \\
5 & Fiber products & Problem Set 5 & Test universal properties and base-change calculations. \\
6 & Properties of morphisms & Problem Set 6 & Test distinctions among finite, separated, affine, and proper morphisms. \\
7 & $\operatorname{Proj}$ & Problem Set 7 & Test graded-algebra/projective-geometry correspondence. \\
8 & Projective schemes & Problem Set 8 & Connect projective constructions with line bundles and embeddings. \\
9 & Quasi-coherent sheaves & Problem Set 9 & Test the module--sheaf correspondence. \\
10 & Coherent sheaves & Problem Set 10 & Test exactness and local properties of coherent sheaves. \\
11 & Derived functors & Problem Set 11 & Test conceptual and technical understanding of sheaf cohomology. \\
12 & Cech cohomology & Problem Set 12 & Test explicit cohomological calculations and affine vanishing. \\
13 & Projective-space cohomology & Problem Set 13 & Test computations with twisting sheaves. \\
14 & Serre vanishing & Problem Set 14 & Test applications of positivity and cohomological vanishing. \\
15 & Synthesis & Problem Set 15 & Integrate the major constructions and proof techniques. \\
Final & Entire course & Oral final examination & Evaluate conceptual mastery, theorem statements, proof strategies, and synthesis. \\
\bottomrule
\end{longtable}

\section{Oral Final Examination Structure}

The oral final should assess the student's ability to organize and communicate the theory rather than merely recall isolated definitions.

A typical examination can contain three components:

\begin{enumerate}
    \item \textbf{Definitions and examples:} The student states and explains several central definitions, such as schemes, morphisms, fiber products, quasi-coherent sheaves, coherent sheaves, and sheaf cohomology, together with examples and counterexamples.
    \item \textbf{Theorems and proof strategies:} The student states major results such as affine vanishing, the correspondence between modules and quasi-coherent sheaves on affine schemes, and Serre vanishing, and explains or reconstructs representative proof ideas.
    \item \textbf{Synthesis problem:} The student is given a geometric situation and asked to determine which scheme-theoretic and cohomological constructions apply, explaining the reasoning step by step.
\end{enumerate}

\section{Global Difficulty Analysis}

\begin{longtable}{p{0.22\textwidth}p{0.35\textwidth}p{0.34\textwidth}}
\toprule
\textbf{Difficult Area} & \textbf{Why Students Struggle} & \textbf{Mitigation} \\
\midrule
Sheaves and stalks & Local and global information are represented simultaneously. & Begin with concrete function sheaves and repeatedly compare sections, stalks, and germs. \\
Affine schemes & Zariski topology and prime ideals behave differently from familiar geometric spaces. & Use explicit spectra such as $\operatorname{Spec}(\mathbb Z)$ and polynomial rings. \\
Schemes & The structure sheaf contains information invisible in the topology. & Compare reduced and nonreduced schemes with identical underlying spaces. \\
Morphisms & A morphism includes both a topological map and a sheaf map. & Track both components in every example. \\
Fiber products & Universal properties and tensor products require simultaneous categorical and algebraic reasoning. & Alternate abstract statements with explicit affine computations. \\
Properties of morphisms & Several definitions are similar but encode different geometric behavior. & Build a running comparison chart with examples and counterexamples. \\
$\operatorname{Proj}$ & Graded localization introduces new notation and algebraic machinery. & Construct projective space explicitly and compare it continuously with $\operatorname{Spec}$. \\
Quasi-coherent sheaves & Students may not see why the definition is natural. & Derive the definition from the affine module--sheaf correspondence. \\
Sheaf cohomology & Derived functors create a substantial abstraction jump. & Motivate cohomology through the failure of global sections to be exact. \\
Cech calculations & Sign conventions and multiple intersections increase cognitive load. & Start with two-open-set covers and expand gradually. \\
Projective cohomology & Computation combines geometry, exact sequences, and combinatorics. & Derive formulas recursively rather than relying on memorization. \\
Serre vanishing & The theorem can appear disconnected from the preceding theory. & Present it as the culmination of positivity, twisting, global generation, and cohomology. \\
\bottomrule
\end{longtable}

\section{Suggested Lecture Sequence}

\begin{longtable}{p{0.08\textwidth}p{0.20\textwidth}p{0.62\textwidth}}
\toprule
\textbf{Week} & \textbf{Lecture} & \textbf{Content} \\
\midrule
1 & 1 & Motivation: from varieties to schemes; presheaves and sheaves. \\
1 & 2 & Stalks, sheafification, and examples of sheaves. \\
2 & 3 & Prime spectrum and the Zariski topology. \\
2 & 4 & Structure sheaf and affine schemes. \\
3 & 5 & Locally ringed spaces and schemes. \\
3 & 6 & Reduced, irreducible, integral, and nonreduced schemes. \\
4 & 7 & Morphisms of schemes. \\
4 & 8 & Closed immersions and affine morphisms. \\
5 & 9 & Products and fiber products. \\
5 & 10 & Fibers and base change. \\
6 & 11 & Finite type, finite, affine, and separated morphisms. \\
6 & 12 & Proper and projective morphisms. \\
7 & 13 & Graded rings and $\operatorname{Proj}$. \\
7 & 14 & Projective space and twisting sheaves. \\
8 & 15 & Projective schemes and closed subschemes. \\
8 & 16 & Quasi-projective schemes and line bundles. \\
9 & 17 & Sheaves of modules and quasi-coherent sheaves. \\
9 & 18 & The affine correspondence between modules and quasi-coherent sheaves. \\
10 & 19 & Coherent sheaves and locally free sheaves. \\
10 & 20 & Exact sequences of sheaves and local exactness. \\
11 & 21 & Global sections and derived functors. \\
11 & 22 & Definition and basic properties of sheaf cohomology. \\
12 & 23 & Cech complexes and cohomology. \\
12 & 24 & Cohomology of quasi-coherent sheaves on affine schemes. \\
13 & 25 & Cohomology of projective space. \\
13 & 26 & Cohomology of twisting sheaves and hyperplane sequences. \\
14 & 27 & Serre vanishing and global generation. \\
14 & 28 & Hilbert polynomials and cohomological applications. \\
15 & 29 & Synthesis: schemes, sheaves, morphisms, and cohomology. \\
15 & 30 & Oral-exam review and integrated examples. \\
\bottomrule
\end{longtable}

\section{Representative Course-Level Questions}

\begin{enumerate}
    \item Explain why the construction $\operatorname{Spec}(A)$ should be regarded as a geometric object associated to the ring $A$, rather than merely as a set of prime ideals.
    
    \item Given a morphism $f:X\to Y$, explain how the induced maps on stalks encode the local algebraic behavior of $f$.
    
    \item Explain the universal property of a fiber product and compute a fiber product of affine schemes using tensor products.
    
    \item Compare affine and projective geometry from the perspective of $\operatorname{Spec}$ and $\operatorname{Proj}$.
    
    \item Explain why quasi-coherent sheaves are the natural sheaf-theoretic analogue of modules on schemes.
    
    \item Why is the global sections functor not exact? Explain how this motivates the definition of sheaf cohomology.
    
    \item State the affine vanishing theorem and explain its geometric significance.
    
    \item Explain how Cech cohomology can make an abstract derived-functor construction computationally accessible.
    
    \item State Serre vanishing and explain the geometric meaning of sufficiently positive twisting.
    
    \item Starting from a coherent sheaf on a projective scheme, describe how one can use twisting, global sections, and cohomology to extract geometric information.
\end{enumerate}

\section{Conceptual Dependency Summary}

The course can be viewed as a sequence of increasingly global layers:

\begin{enumerate}
    \item \textbf{Local algebra:} rings, localization, and prime ideals.
    \item \textbf{Local geometry:} $\operatorname{Spec}(A)$ and locally ringed spaces.
    \item \textbf{Global geometry:} schemes and morphisms.
    \item \textbf{Geometric constructions:} fiber products, base change, and projective schemes.
    \item \textbf{Geometric algebra:} quasi-coherent and coherent sheaves.
    \item \textbf{Global invariants:} sheaf cohomology.
    \item \textbf{Applications:} projective-space cohomology, Serre vanishing, and Hilbert-polynomial methods.
\end{enumerate}

The central conceptual transition is from the classical viewpoint
\[
\text{varieties} \longrightarrow \text{schemes}
\]
followed by
\[
\text{regular functions} \longrightarrow \text{sheaves}
\]
and ultimately
\[
\text{global sections} \longrightarrow \text{cohomology}.
\]

This progression provides the organizing principle for the course: schemes enlarge the class of geometric objects, sheaves encode local algebraic data, and cohomology measures the global obstructions that remain after local information has been understood.

\end{document}